\subsection{自定义Poisson目标的R扩展路径}
安装版示例以二维Poisson回归为例，$y_i\sim\operatorname{Poisson}(\exp(X_i\beta))$、$\beta_j\sim N(0,4)$。外部Python模型提供框架密度、独立NumPy密度/解析梯度和恒等参数变换；R整批调用，再转换为\texttt{posterior}对象。以下为仓库中实际示例的调用顺序，环境变量在启动R前指定，输出使用不存在的新目录。
\begingroup\small
\begin{verbatim}
# Once: create the shared data and actual random arrays
python examples/installed_custom_target.py prepare \
  --output output/poisson-inputs
# Set R_LIBS_USER and RETICULATE_PYTHON before starting R
Rscript --vanilla examples/installed-custom-target.R \
  torch output/poisson-inputs output/poisson-fit cpu
\end{verbatim}
\endgroup
以下片段展示与上述R命令等价的调用，两者选择其一运行。从仓库根目录出发，目录与返回值的关系如下。核和执行器在Python示例的\texttt{workflows}配置中选择，依次为MALA顺序/quasi-DEER、RWM顺序/Picard；主示例对四种组合全部运行。
\begingroup\small
\begin{verbatim}
example_dir <- normalizePath("examples", mustWork = TRUE)
inputs <- normalizePath("output/poisson-inputs", mustWork = TRUE)
output <- "output/poisson-fit" # Must not already exist
library(parallelbayes)
pb_environment("torch")
helper <- reticulate::import_from_path(
  "installed_custom_target", path = example_dir, convert = TRUE)
result <- helper$run_example(inputs, output, "torch", "cpu")
print(vapply(result$workflows, function(w) w$status, ""))
for (w in result$workflows) {
  if (!identical(w$status, "completed")) next
  a <- w$draws
  dimnames(a) <- list(NULL, NULL, unlist(w$names))
  draws <- posterior::as_draws_array(a)
  print(dim(draws)) # iteration, chain, parameter
  print(posterior::summarise_draws(draws))
}
\end{verbatim}
\endgroup
实际Mac torch安装示例返回四个\texttt{completed}状态，每组样本维度为\texttt{128 × 4 × 2}。第一组顺序MALA的\texttt{beta0}均值为0.11428、标准差0.16642、$\hat R=1.01917$，\texttt{beta1}的$\hat R=1.03180$，说明接口成功返回仍可伴随短链诊断问题。完整脚本同时核对安装模块来源和输入哈希，保存失败项、样本及诊断，并在数值检查失败时返回错误。该例已在Mac安装环境下运行JAX/torch与R入口，四链各128次转移用于说明接口；没有预热，不能作为推荐推断预算。它的Windows/CUDA组合尚未实测。水井案例另有Windows CPU/CUDA/R接口证据。\texttt{pb\_model()}不接收任意R闭包，Pyro CPU NUTS也仍是独立研究CLI。
